\bookchapter{第十七章：手机与其他设备的联动}

\chapterintro{你手里的手机，其实是一个超级遥控器兼数据中心。它不光能打电话，还能把电脑、电视、手表、音箱、汽车串在一起，让这些设备不再是孤岛。这章我们就来聊聊怎么让它们“组队打怪”。}

\booksection{一、与电脑协同：手机和电脑不再是两个世界}

说实话，我以前最烦的就是手机拍了个照片，想放到电脑上修一下，还得找数据线、等传输。后来发现，其实根本不需要数据线，好的联动比数据线快十倍。

\begin{itemize}
  \item \textbf{苹果生态（AirDrop + 通用剪贴板）}：如果你用iPhone配Mac，隔空投送是你第一个该学会的功能。在手机上点分享 隔空投送 选你的Mac，照片瞬间就过去了，不用流量不用Wi-Fi，用的是蓝牙和Wi-Fi直连。还有个容易被忽略的功能叫"通用剪贴板"——手机上复制一段文字，在Mac上直接粘贴就出来了，简直像魔法一样。
  \item \textbf{安卓+Windows（连接至Windows）}：微软有个叫"连接至Windows"的App，三星手机内置，其他安卓手机去应用商店下一个就行。配对之后你可以在电脑上直接看手机通知、回短信、翻照片，甚至接打电话。我自己用了两年，开会的时候不用一直低头看手机，电脑屏幕上弹个通知就够了。
  \item \textbf{品牌多屏协同}：华为的"多屏协同"和小米的"小米妙享"更狠——直接把手机屏幕投到电脑上，你可以用鼠标拖拽文件从手机到电脑，也能在电脑上刷抖音、回微信。跨设备复制粘贴、文件拖拽都行，体验非常丝滑。
\end{itemize}

\booksection{二、投屏到电视：全家一起看}

过年回家的时候，我堂妹想在电视上放她手机里的毕业旅行视频。折腾了半天，最后发现其实非常简单：

\begin{itemize}
  \item \textbf{iPhone 电视}：从右上角下拉控制中心，点那个两个方框叠在一起的图标（屏幕镜像），找到你的电视名称，一点就连上了。前提是电视和手机连的同一个Wi-Fi，而且电视支持AirPlay（大部分智能电视都支持）。
  \item \textbf{安卓 电视}：各家叫法不同，华为叫"无线投屏"，小米叫"投射"，OPPO/vivo叫"多屏互动"。一般在快捷开关面板里能找到，找不到就去设置里搜"投屏"。
  \item \textbf{视频App直接投}：如果你只是想在电视上看爱奇艺、B站、YouTube，其实更简单——在播放界面找一个矩形带Wi-Fi符号的小图标，点一下，选电视，视频就过去了。这时候手机可以干别的，视频在电视上继续放，互不影响。
  \item \textbf{投屏不成功怎么办}：90\%的情况是没连同一个Wi-Fi。还有10\%是电视的投屏功能没开，去电视设置里找"多屏互动"或"AirPlay"开关。
\end{itemize}

\booksection{三、与智能手表/手环：健康管家长在手腕上}

我给我妈买了个手环，她一开始说“我又不运动买这个干嘛”。结果用了两周，她发现自己每天只睡五个多小时，开始主动调整作息了。这就是手表/手环最实用的价值——把你看不见的健康数据变成看得见的数字。

\begin{itemize}
  \item \textbf{通知过滤器}：手机在包里或口袋里震个不停的时候，抬手腕看一眼——是老板的微信还是拼多多的促销，一秒判断，不值得掏手机的消息就让它去。
  \item \textbf{健康数据同步}：手表24小时监测心率、血氧、睡眠、步数，这些数据会自动汇总到手机上的健康App里。久而久之你能看到自己的趋势——这个月是不是走得太少，这周睡眠质量有没有变好。
  \item \textbf{交通卡/门禁卡同步}：在手机上绑好的公交卡和门禁卡，大部分品牌会自动同步到手表上。坐地铁的时候不用掏手机，手腕一靠闸机就过去了，真的会有种"未来已来"的感觉。
  \item \textbf{找手机神器}：手表上一般都有"找手机"功能，点一下手机就会响。我自己用这个功能的频率远超预期——手机随手一放就忘在哪了。
\end{itemize}

\booksection{四、与智能家居：手机是家里的“总开关”}

智能家居听起来很玄乎，但说穿了就三件事：用手机远程控制、设定自动化规则、语音操控。

说说我自己的几个实际用法，你就知道这玩意儿多实用了：

\begin{itemize}
  \item \textbf{空调}：夏天回家路上提前用手机把空调打开，进门就是凉快的。花一百多买个空调伴侣插座就能实现，不用换空调。
  \item \textbf{灯}：睡前躺床上，懒得起来关灯，手机上点一下或者喊一声"嘿Siri/小爱同学，关灯"。
  \item \textbf{扫地机器人}：出门前手机上点一下"开始清扫"，回来地就干净了。还能设每天定时扫。
  \item \textbf{摄像头}：家里有老人小孩或者宠物的，手机上随时能看一眼家里的情况。
\end{itemize}

选哪个生态？一句话总结：

\begin{itemize}
  \item 用iPhone 优先选支持Apple HomeKit的设备（贵一点但隐私保护好）。
  \item 设备多求性价比 小米米家生态最全，灯、插座、扫地机什么都有。
  \item 用华为手机 优先鸿蒙智联，和华为设备体验最好。
\end{itemize}

\booksection{五、蓝牙耳机与音箱：无线的自由}

从有线耳机换成蓝牙耳机，对我来说是手机体验升级里最明显的一个。再也不会从包里掏出一团乱麻，也不用担心线挂在门把手上了。

\begin{itemize}
  \item \textbf{配对很简单}：耳机第一次开机一般自动进入配对模式（闪灯），手机打开蓝牙就能搜到，点一下就连上了。之后每次拿出来自动连，什么都不用做。
  \item \textbf{同品牌有加成}：AirPods配iPhone会自动弹窗显示电量，还能在多个苹果设备之间无缝切换。华为FreeBuds配华为手机、小米耳机配小米手机也都有类似的弹窗和优化。
  \item \textbf{蓝牙音箱}：家庭聚会、户外露营的时候，一个蓝牙音箱能让气氛翻倍。连上之后，手机就是你的DJ台。
  \item \textbf{注意延迟}：打游戏的时候如果用蓝牙耳机，可能会有轻微延迟（枪响了半秒才听到枪声）。如果对延迟敏感，要么用有线耳机，要么买支持"低延迟模式"的蓝牙耳机。
\end{itemize}

\begin{tcbinfo}{UWB超宽带技术——比蓝牙更精准的"空间感"}
除了蓝牙，新款手机里还藏着一颗UWB（Ultra-Wideband，超宽带）芯片。它的能力是\textbf{精确测距和空间定位}——能算出另一个设备离你多远、在哪个方向，精度可达厘米级。蓝牙只能告诉你“在附近”，UWB能告诉你“在你左前方1.2米”。常见应用：

\begin{itemize}
  \item \textbf{AirTag/查找精确}：iPhone查找AirTag时显示"7米远、向左走"的箭头导航，靠的就是UWB。普通蓝牙只能告诉"在不在范围内"。
  \item \textbf{数字车钥匙}：搭载UWB的iPhone和安卓旗舰，可以当车钥匙走近车辆自动解锁、走远自动锁车，还能区分驾驶员和乘客位置，防中继破解。
  \item \textbf{室内导航}：商场、机场里通过UWB信标实现"过道级"导航，比GPS准得多。
  \item \textbf{支持的机型}：iPhone 11及以后机型均带UWB芯片（U1/U2）；安卓阵营小米、三星、vivo、OPPO部分旗舰也已搭载。
\end{itemize}

\end{tcbinfo}

\begin{tcbinfo}{LE Audio与Auracast——蓝牙音频的新一代标准}
用了二十年的经典蓝牙音频（A2DP）正在被\textbf{LE Audio}（低功耗音频）取代，这是蓝牙5.2以后的新一代音频标准：

\begin{itemize}
  \item \textbf{更省电、音质更好}：采用全新的LC3编码，比老标准的SBC音质更好、功耗更低，耳机续航明显延长。
  \item \textbf{双路音频}：左右耳塞独立传输，单耳模式切换不掉线。
  \item \textbf{Auracast广播音频}：LE Audio的招牌功能。一个手机/电视可以同时向\textbf{无数副}耳机广播音频——机场候机厅、博物馆、健身房、教室里都能用一个发射源，多人用自己的耳机听同一画面。
  \item \textbf{现状}：2024-2026年新上市的旗舰耳机和手机普遍支持LE Audio，老设备需硬件升级无法软件解锁。买新耳机时可以留意是否支持LE Audio/Auracast标识。
\end{itemize}

\end{tcbinfo}

\booksection{六、与车载系统：开车时的安全搭档}

开车看手机是违法的，也是要命的。但有时候确实需要导航、接电话、放音乐。这时候CarPlay（苹果）和Android Auto（安卓）就派上用场了。

\begin{itemize}
  \item \textbf{怎么连}：如果你的车支持（近几年大部分新车都支持），用数据线把手机插到车上的USB口，车载屏幕会自动切换成CarPlay或Android Auto界面。部分新车支持无线连接，更省事。
  \item \textbf{能干嘛}：导航（比车载导航好用多了）、打电话（方向盘接听）、听歌（手机里的音乐App直接投到车上）、语音发微信（Siri或小爱帮忙读消息和回复）。
  \item \textbf{老车也能用}：如果你的车比较老不支持，花几百块买个外挂的CarPlay屏幕，粘在仪表台上，体验差不多。
\end{itemize}

\booksection{七、联动中的安全和隐私}

设备联的越多，隐私泄露的风险面就越大。几点需要注意：

\begin{itemize}
  \item 公共场合投屏之前，先看看手机通知栏有没有私密消息，别让全家看到不该看的内容。
  \item 把设备送给别人或者卖掉之前，先解绑——去手机设置里找到那个设备，选"移除"或"取消配对"。
  \item 智能家居摄像头，买有品牌的、有加密的，别买几十块的杂牌。听说过太多摄像头被黑客直播的新闻了。
\end{itemize}

\begin{tcbpractice}{来试试：把手机屏幕投到电视上，并设置一个智能场景}
\begin{enumerate}
  \item 确保手机和电视连同一个Wi-Fi。
  \item 在电视上找到"屏幕镜像"或"AirPlay"功能并打开。
  \item 在手机上：iPhone从控制中心点「屏幕镜像」；安卓下拉找「无线投屏」。
  \item 选一张你最喜欢的照片投到电视上，让家人看看。
  \item 如果有智能家居设备，试着在手机App里创建一个"离家模式"：一键关灯、关空调、启动扫地机器人。
\end{enumerate}

\end{tcbpractice}

\begin{tcbinfo}{案例 1：电池续航优化实战}
\begin{itemize}
  \item \textbf{场景}：用户小米14 续航差，半天就没电。
  \item \textbf{步骤}：(1) 设置 电池 查看耗电排行 (2) 识别异常后台 App (3) 关闭 5G 切 4G (4) 限制后台 (5) 开启省电模式 (6) 充电时摘壳散热。
  \item \textbf{预期结果}：续航延长 30-50\%，电池健康度数据可读。
\end{itemize}

\end{tcbinfo}
